\begin{figure}[htbp]
\centering
\begin{tikzpicture}[>=Latex,x=1cm,y=1cm,every node/.style={font=\small},dot/.style={circle,fill,inner sep=1.6pt}]
\node[dot,label=above:$u$] (u) at (0,2.6) {};
\node[dot,label=left:$v$] (v) at (-.9,1.3) {};
\coordinate (a) at (-1.7,0);
\coordinate (b) at (1.3,0);
\coordinate (q) at (-.1,0);
\draw[thick] (u)--(v)--(a) node[below] {$a$};
\draw[thick] (u)--(b) node[below] {$b$};
\draw[dashed] (v)--(q) node[below,align=center] {extra\\branch};
\node at (0,-1.15) {Before marginalization};
\draw[->] (2.1,1.4)--(3.7,1.4) node[midway,above,align=center,font=\footnotesize] {replace aliases;\\prune};
\node[dot,label=above:$u$] (u2) at (5.5,2.6) {};
\coordinate (a2) at (3.8,0);
\coordinate (b2) at (6.8,0);
\draw[thick] (u2)--(a2) node[below] {$a$};
\draw[thick] (u2)--(b2) node[below] {$b$};
\node[draw,circle,inner sep=2pt,label=left:$v$] at ($(u2)!.5!(a2)$) {};
\node at (5.5,-1.15) {The retained law has the fixed split $u$};
\end{tikzpicture}
\caption{The earlier-pair configuration in the centered integration argument. The pair $(a,b)$ splits at the fixed mass $u<v$. Replacing the aliases and integrating out the extra branch leaves the two paths on the right. Their joint law is independent of $v$, which now marks an evaluation point rather than a retained fork. Values at fixed projection cuts remain observables of these paths. The single-path estimate is applied under this fixed law before H\"older's inequality handles the other random factors.}
\label{fig:fixed-marginal}
\end{figure}
